\documentclass[runningheads]{llncs}

\emergencystretch 1.5em % tolerance for overfull lines; remove for strict typesetting

\usepackage[T1]{fontenc}
\usepackage{graphicx}
\usepackage{epsfig}
\usepackage{subcaption}
\usepackage{calc}
\usepackage{amsmath}
\usepackage{amssymb}
\usepackage{amstext}
\usepackage{multicol}
%\usepackage{pslatex} % disabled 2026-09-16: Symbol-font bold + \vec accent warnings; default LLNCS (CM) fonts used instead. Re-enable if Times look is wanted.
\usepackage{algorithm2e}
\usepackage[bottom]{footmisc}
\usepackage{booktabs}
\usepackage{multirow}
\usepackage{svg}
\usepackage{hyperref}

\pagenumbering{gobble}
\begin{document}
%
\title{PAMR: Link Prediction on Knowledge Graphs via Path-Aware Multi-modal Reasoning}
%
\titlerunning{Link Prediction on KG via Path-Aware Multi-modal Reasoning}

\author{Hai-Dang Huynh-Lam\inst{1,2}\orcidID{0000-0001-9487-7281} \and
  Ngoc-Thao Nguyen\inst{1,2} \and
Ngoc-Thanh Le\inst{1,2}\orcidID{0009-0009-6426-2652}}
%
\authorrunning{Hai-Dang et al.}
\institute{
  Faculty of Information Technology, University of Science, Ho Chi Minh City, Vietnam
  \and
  Vietnam National University, Ho Chi Minh City, Vietnam\\
\email{\{hlhdang,nnthao,lnthanh\}@fit.hcmus.edu.vn}}
%
\maketitle  % typeset the header of the contribution
%
\begin{abstract}
  Knowledge graphs organize facts as interconnected entities and relations, yet
  they are always incomplete. Multimodal knowledge graphs fill the gaps with
  images and text, but prevailing link prediction models fold these signals into
  a single vector per entity, which hides the connectivity that carries meaning
  across modalities. We ask what happens when the attachments themselves become
  graph nodes. Under this modality-as-entity formulation, our framework,
  Path-Aware Multimodal Reasoning (PAMR), extracts shortest paths between query
  entities and scores each path with a Transformer encoder over frozen knowledge
  graph embedding tokens. Separating path modelling from embedding training
  avoids the initialization cost of complex multimodal representations. On
  WikiArt-MKG v2, over the four cross-modality artist relations that the dataset
  offers, PAMR reaches an MRR of 0.379 (a 4.4\% gain over the strongest embedding
  baseline) and the best Hits@5 at 0.782. Each prediction comes with the
  multimodal path that supported it, so every prediction stays traceable to
  actual artworks, artists, and shared attributes.

  \keywords{Knowledge Graph \and Multimodal Learning \and Path-based Reasoning \and Link Prediction \and Transformers.}
\end{abstract}
%
%
%

\section{Introduction}
\label{sec:introduction}

Knowledge graphs (KGs) represent facts as entities wired together by typed relations, and link prediction asks for the connections that are missing. Most methods learn continuous representations from structure alone. Multimodal knowledge graphs (MMKGs) add images and text, and multimodal embedding models do lift link prediction accuracy, but they fold visual and textual signals into one vector per entity. The folding hides connectivity patterns that cross modality boundaries and leaves reasoning hard to inspect.

Path-aware reasoning goes the opposite way: it reads explicit structural sequences between two entities rather than a single proximity score. Yet the models built for this are designed around concept-only graphs, and they tend to tie path modelling to embedding training, a coupling that becomes expensive once multimodal representations are involved. Neither limitation is forced by the data: the graph could instead carry multimodal values as nodes.

We therefore propose the Path-Aware Multimodal Reasoning framework (PAMR). PAMR solves link prediction on multimodal knowledge graphs by scoring structural paths between query entities with a Transformer encoder whose tokens are frozen knowledge graph embeddings. Under a modality-as-entity setting, multimodal information becomes part of the topology, so extracted paths can route through shared artworks, styles, or other attributes. The evidence is explicit and traceable, and the decoder-free design sidesteps the initialization deadlock of continuous sequence modelling.

The core contributions of this work are summarized as follows.

\begin{itemize}
  \item First, we present an architectural framework for path-aware reasoning for link prediction on MMKGs under a modality as entity setting, establishing a robust baseline utilizing the unified embeddings provided by VL-KGE.
  \item Second, the proposed path extraction uses both conceptual entities and multimodal entities, with experimental results demonstrating substantial performance gains derived from integrating these multimodal signals.
  \item Third, we introduce a Transformer based sequential architecture for scoring the extracted paths. This approach achieves results comparable to strong embedding baselines while complementing each prediction with explicit, traceable multimodal reasoning paths, even though the learned Transformer scorer itself is not inherently transparent.
\end{itemize}

\section{Related Work}
\label{sec:related-work}

The PAMR framework sits at the intersection of graph-based relational reasoning, multimodal representation learning, and sequential modelling. We survey each line and state the gap PAMR fills.

\subsection{Path-Based Reasoning on Knowledge Graphs}
\label{sec:rw-path-reasoning}

Path-based reasoning extracts and evaluates multi-hop relational sequences between entities instead of predicting single edges. The lineage starts with statistical random walks such as the Path Ranking Algorithm~\cite{lao2011random}, moves through differentiable logic (NeuralLP~\cite{yang2017differentiable}, RNNLogic~\cite{qu2021rnnlogic}), and reaches reinforcement learning navigators where DeepPath~\cite{xiong2017deeppath} and MINERVA~\cite{das2018go} cast pathfinding as a Markov decision process. Variational inference~\cite{chen2018variational} and hierarchical multi-agent RL~\cite{zhang2022learning} later answered the sparse-reward problem of long paths.

GNN- and sequence-based encoders followed. RED-GNN and NBFNet parameterize the Bellman-Ford recursion with neural components~\cite{zhang2022knowledge,zhu2021neural}, AdaProp prunes the propagation~\cite{zhang2023adaprop}, and SQUIRE translates queries directly into paths with a Transformer~\cite{bai2022squire}. Across all of these, two habits persist: embeddings stay learned rather than frozen, and scoring is driven by message passing, generative decoding, or RL rather than path-sequence Transformers. And no method is built for multimodal entities, which keeps the embedding-sequence entanglement costly when richer representations enter.

\subsection{Multimodal Knowledge Graph Embeddings}
\label{sec:rw-mmkge}

Multimodal knowledge graph embedding (MMKGE) targets the structural sparsity of KGs by injecting images and text into representation learning. IKRL~\cite{xie2017image} and TransAE~\cite{wang2019transae} first aligned visual and textual features with structural embeddings, and the MMKG collection~\cite{liu2019mmkg} supplied matching baselines. Later systems manage modality noise and structural complexity: RSME forgets conflicting signals with a gate~\cite{wang2021rsme}, OTE preserves topology through orthogonal relation transforms~\cite{tang2020orthogonal}, and MMRNS biases negative sampling with cross-modal attention~\cite{xu2022multimodal}.

Transformer-era MMKGE pushes fusion further. MKGformer fuses visual and textual views at multiple levels~\cite{chen2022mkgformer}, MoSE splits modalities before joining them~\cite{zhao2022mose}, VISTA encodes at the triplet level with decoding Transformers~\cite{lee2023vista}, and VL-KGE turns to vision-language models for modality asymmetry at scale~\cite{efthymiou2026vlkge}. One property survives in all of them: modalities are fused into a single entity embedding that serves triple scoring, not path discrimination. No public benchmark we know of treats multimodal values as entities that connect several conceptual nodes at once, and that is exactly the setting PAMR evaluates on.

\subsection{Transformers for Knowledge Graph Reasoning}
\label{sec:rw-transformers}

Transformers and pre-trained language models let structural data be contextualised through self-attention. KG-BERT reads triples as text for sequence classification~\cite{yao2019kg}, and StAR splits contexts to cheapen cross-encoding every candidate~\cite{wang2021star}. CoKE approaches the graph natively, tokenizing edges and paths into encoder input~\cite{wang2019coke}, while HittER aggregates neighbourhoods hierarchically through masked entity prediction~\cite{chen2021hitter}.

For complex logical queries, KG-Transformer pairs masked pretraining with a sparse Mixture-of-Experts activation~\cite{liu2022mask}. Generative formulations round out the picture: KGT5 casts link prediction as sequence-to-sequence~\cite{saxena2022kgt5}, GenKGC generates targets from relation-guided prompts~\cite{zhu2022genkgc}, and CSPrompter bridges symbolic logic and large language models through structure prompting~\cite{chen2023csprompter}. What these systems encode are triples, subgraphs, or text, however, not interleaved multimodal path sequences. They open with language-model pretraining rather than frozen KGE tokens, and they aim at text-only graphs. The multimodal-as-entity graphs PAMR targets fall outside this scope as well.

% \subsection{Evaluation Approaches}
% \label{sec:rw-evaluation}
% [Commented draft of negative-sampling survey removed for the SOICT version;
%  see assets/paper.tex for the original text.]

Taken together, the three lines leave open questions that PAMR takes up: how to run path-aware reasoning over heterogeneous multimodal entities, how to project frozen KGE embeddings into a Transformer's token space without an initialization dead end, and how to keep sequence ranking stable with a perfectly balanced, query-level paired sampling scheme.

\section{Method}
\label{sec:method}
\subsection{Overview}
\label{sec:overview}

A KG is a directed multi-relational graph
$\mathcal{G} = (\mathcal{V}, \mathcal{E}, \mathcal{R})$ where
$\mathcal{V}$ is a set of entities, $\mathcal{R}$ a set of relation types,
and $\mathcal{E} \subseteq \mathcal{V} \times \mathcal{R} \times \mathcal{V}$
a set of labelled edges. Each edge $(s, r, o) \in \mathcal{E}$
denotes that subject entity $s$ is connected to object entity $o$
via relation $r$.

Path-based reasoning focuses on using structural paths and connection
patterns within KGs to infer new relations or answer queries. Given
a query $(s, r, ?)$, the target is to approximate the conditional
distribution $P(\hat{o} \mid s, r, \mathcal{G})$ over candidate
objects $\hat{\mathbb{O}}$. By using the shortest paths $p_o$ between $s$ and $\hat{o}$ in $\mathcal{G}$ as the most important factor contributing to the distribution,
the target reduces to $P(p_o \mid s, r)$. A scoring function
$\hat{f}_\theta$ is learned to predict this target and then assigns a real-valued score to each path given the
query relation $r$, afterward the candidate with the highest-scoring path is selected as the answer.

In a multimodal KG the entities differ in nature; artworks are inherently visual, while artists, styles, genres, and locations are textual. The prevailing design, multimodal-as-attachment, keeps one node per entity and aggregates modality-specific embeddings $\mathbf{v}_i$, indexed by $i \in \mathbb{M}$, into a single point $\mathbf{v} \in \mathbb{H}$ (Section~\ref{sec:embedding-module} gives the aggregation). We instead split every multimodal value into its own node, the modality-as-entity formulation, so the topology itself carries the modality structure; a mapping $\chi: \mathcal{V} \to \mathbb{M}$ records which modality produced each node.

This construction is formalized compactly as follows. Each entity carries an entity type $\tau: \mathcal{V} \to \mathcal{T}$ (e.g., artwork, artist, style, genre, year, location), and a modality label $\mu(v) \in \mathbb{M}$ defined only on modality-bearing types $\mathcal{T}_{\text{mod}} \subseteq \mathcal{T}$. The modality-as-entity transformation externalizes every modality-bearing attribute value into a first-class node $\tilde{v}$ with a single modality label $\mu(\tilde{v})$, yielding a transformed graph $\mathcal{G}^{*} = (\mathcal{V}^{*}, \mathcal{E}^{*}, \mathcal{R}^{*})$ in which references to attribute values are mediated through these nodes. The transformation obeys four structural constraints: (i) each modality node carries exactly one modality label $\mu(\tilde{v})$; (ii) modality nodes interface with conceptual entities exclusively through dedicated mediation relations (e.g., \texttt{hasArtist} links an artwork node to artist nodes); (iii) conceptual relations connect conceptual entities only, so path traversal between conceptual entities may cross type boundaries solely via modality-mediated hops; and (iv) a single modality node may mediate several conceptual entities at once, a sharing that per-entity attachment cannot express and that turns multimodal values into first-class bridges between concepts.

The proposed method, PAMR, consists of four components as illustrated in Figure \ref{fig:overview}. First, a frozen
pre-trained KGE model supplies token embeddings for entities and relations.
Second, a bidirectional BFS extracts shortest paths between query pairs under the modality-as-entity setting.
Third, a Transformer encoder scores each path by contextualising
interleaved entity and relation tokens. Fourth, a margin-based ranking
loss trains the scorer on paired positive and negative paths. Each
component is detailed in the following subsections.

\begin{figure}[t]
  \centering
  \includesvg[width=\linewidth]{figures/overview.svg}
  \caption{Overview of PAMR. Given a query $(s, r, ?)$, frozen KGE embeddings are retrieved for entities and relations. A bidirectional BFS extracts shortest paths between the subject and each candidate object under the modality-as-entity setting. Each path is encoded as an interleaved sequence of entity and relation tokens, prepended with a CLS token and the query relation. The Transformer encoder contextualises the sequence, and a two-layer MLP readout produces a scalar score per path. Paired positive and negative paths are scored jointly and trained with a margin-based ranking loss.}
  \label{fig:overview}
\end{figure}

\subsection{Embedding Module}
\label{sec:embedding-module}

To supply pre-encoded relational structure without incurring additional
training cost, we use a frozen pre-trained KGE model.  A KGE model
learns entity embeddings $\mathbf{E} \in \mathbb{H}^{|\mathcal{V}|}$
and relation embeddings $\mathbf{R} \in \mathbb{H}^{|\mathcal{R}|}$
such that a scoring function $g: \mathcal{V} \times \mathcal{R} \times
\mathcal{V} \to \mathbb{R}$ assigns higher values to true triples
than to false ones.  The embedding space $\mathbb{H}$ is a
$D$-dimensional continuous manifold.

In a multimodal KG, entities possess heterogeneous modalities.
VL-KGE~\cite{efthymiou2026vlkge} constructs a unified embedding for
every entity by fusing its available modalities into a single point
in $\mathbb{H}$.  For each entity $v \in \mathcal{V}$, let
$\mathcal{M}(v) \subseteq \mathbb{M}$ denote the set of available
modalities.  The unified entity embedding $\mathbf{v}$ is obtained
by averaging over available modalities:
$\mathbf{v} = \frac{1}{|\mathcal{M}(v)|} \sum_{i \in \mathcal{M}(v)} \mathbf{v}_i$,
where $\mathbf{v}_i$ is the modality-specific embedding for modality
$i \in \mathbb{M}$.  This per-entity averaging naturally handles
modality asymmetry, where different entity types carry different
modality combinations without requiring a fixed global fusion scheme.

The effective dimensionality of entity and relation representations
depends on the KGE architecture: complex-valued models output $2D$
dimensions while real-valued models output $D$. The Transformer
hidden dimension is set to $d = \max(d_{\mathcal{V}},
d_{\mathcal{R}})$, and two learnable linear projections align entity
and relation token dimensions to $d$. These
projections are learned during PAMR training while the KGE model
remains frozen.

\subsection{Path Extraction}
\label{sec:path-extraction}

To provide the Transformer with explicit, interpretable connectivity evidence rather than
relying solely on the implicit proximity captured by KGE embeddings, we extract shortest
paths between query pairs via bidirectional breadth-first search (BFS), bounded by a
maximum hop count $K$.

\paragraph{Multimodal and non-multimodal configurations.}
We define two graph configurations that produce qualitatively different path spaces.
The multimodal configuration $\mathcal{G}_{\text{multi}}$ comprises the complete
graph with all entity types, where paths may bridge across different entity types,
exploiting cross-type relationships for reasoning.  The concept-only configuration
$\mathcal{G}_{\text{concept}}$ restricts the graph to conceptual entities and
edges where both endpoints are conceptual, confining paths to a homogeneous
subgraph that excludes multimodal information.  The latter corresponds to
Path-Aware Reasoning (PAR), which relies solely on structural connectivity
without multimodal signals.

Notably, under the modality-as-entity formulation, multimodal signals and
graph topology are entangled by construction: multimodal information exists
only as explicit nodes, so any configuration that removes or masks multimodal
token features necessarily also removes the corresponding nodes and their
incident edges. A controlled ablation that retains the full topology while
eliminating multimodal features is therefore not defined in this setting.
The concept-only configuration (PAR) serves precisely as this topology-only
counterpart, in which multimodal nodes and their edges are jointly removed,
isolating the contribution of modality-mediated connectivity. The one
alternative that preserves connectivity, tokenizing modality nodes with a
constant placeholder, turns every such node into an information-less dummy
entity, because a modality node's meaning lies precisely in which concepts
it mediates; such an ablation would measure a corrupted graph rather than
the absence of multimodal features, and PAR remains the coherent control
in this formulation.

Formally, a path $p$ between $s$ and $\hat{o}$ comprises $\ell + 1$
entities connected by $\ell$ edges, where
$\ell$ denotes the hop count bounded by $K$.
Paths with $\ell \leq 0$ are excluded from training and evaluation as they
carry no structural reasoning signal, which include unreachable pairs and
trivial self loops.

\subsection{PAMR Architecture}
\label{sec:pamr-architecture}

\subsubsection{Token Interleaving}
\label{sec:token-interleaving}

To encode path structure as a flat sequence, entity and relation tokens are interleaved
with a classification token and the query relation prepended.
Let $\mathbf{v}_j \in \mathbb{R}^{d}$ be the entity token for entity $j$ in the path,
obtained from the unified entity embedding via a learnable projection.
A path with $\ell$ hops comprises $\ell+1$ entity tokens and $\ell$ edge tokens.
These are interleaved into a token sequence of length $L = 2\ell + 3$:
\begin{equation}
  \mathbf{T} = \bigl[\mathbf{c};\; \mathbf{r};\;
    \mathbf{v}_{s};\; \mathbf{r}_{1};\; \mathbf{v}_{2};\; \mathbf{r}_{2};\;
  \mathbf{v}_{3};\; \ldots;\; \mathbf{v}_{o}\bigr]
  \in \mathbb{R}^{L \times d}
  \label{eq:token-seq}
\end{equation}
Here $\mathbf{c} \in \mathbb{R}^{d}$ is a learnable classification token
at position 0, initialised from $\mathcal{N}(0, 0.02^2)$.
$\mathbf{r}$ is the query relation representation at position 1, while entity tokens
$\mathbf{v}_{s}, \mathbf{v}_{2}, \ldots, \mathbf{v}_{o}$ occupy positions $2 + 2i$
for $i = 0, \ldots, \ell$, and edge tokens $\mathbf{r}_{1}, \ldots, \mathbf{r}_{\ell}$
occupy positions $3 + 2i$ for $i = 0, \ldots, \ell - 1$.
All positions receive absolute positional embeddings
$\mathbf{p}_i = \operatorname{PE}(i)$ via a learned embedding
$\operatorname{PE}: \{0, \ldots, 2K+2\} \to \mathbb{R}^{d}$,
where $K$ is the maximum hop count.
The final token input to the Transformer is $\mathbf{T}^{\text{in}}_i = \mathbf{T}_i \;+\; \mathbf{p}_i$.
\subsubsection{Transformer Encoder}

\label{sec:transformer-encoder}

The token sequence is contextualised through a stack of $N$
Transformer encoder layers~\cite{vaswani2017attention} with $h$ attention heads each:
\begin{equation}
  \mathbf{H}^{(0)} = \mathbf{T}^{\text{in}},\qquad
  \mathbf{H}^{(l)} = \operatorname{EncLayer}^{(l)}\bigl(\mathbf{H}^{(l-1)},\; \mathbf{M}\bigr)
  \quad (l = 1, \ldots, N)
  \label{eq:transformer}
\end{equation}
Here $\mathbf{M} \in \{0, 1\}^{L}$ is a padding mask ($\mathbf{M}_i = 1$ indicates an
invalid position) for paths shorter than the batch maximum $N_{\max}(\mathcal{B})$.
Each encoder layer applies:
\begin{align}
  \tilde{\mathbf{H}} &= \operatorname{LayerNorm}\bigl(
  \mathbf{H}^{(l-1)} + \operatorname{MHA}(\mathbf{H}^{(l-1)}; \mathbf{M})\bigr)
  \label{eq:attn}\\
  \mathbf{H}^{(l)} &= \operatorname{LayerNorm}\bigl(
  \tilde{\mathbf{H}} + \operatorname{FFN}(\tilde{\mathbf{H}})\bigr)
  \label{eq:ffn}
\end{align}
with $\operatorname{FFN}(\mathbf{x}) = \mathbf{W}^{(2)} \operatorname{GELU}(\mathbf{W}^{(1)} \mathbf{x} + \mathbf{b}^{(1)}) + \mathbf{b}^{(2)}$
and inner dimension $d'$, set to $2048$ in all runs.
GELU \cite{hendrycks2016gaussian} activation
and Post-LN layer normalisation are used throughout.

\subsubsection{Readout and Scoring}

\label{sec:readout}

To obtain a fixed-dimensional path representation, the contextualised token
sequence is aggregated from
$\mathbf{H}^{(N)} \in \mathbb{R}^{L \times d}$
via one of two readout strategies: CLS extraction, which takes the
classification token at position 0, and mean pooling, which averages
over valid tokens.
A two-layer MLP maps the aggregated representation to a scalar score
$\hat{y} = \mathbf{w}_2^{\top} \operatorname{GELU}(\mathbf{W}_1 \mathbf{z} + \mathbf{b}_1) + b_2$
with dropout $p = 0.1$ applied after GELU.
All linear layers use Xavier
uniform initialisation and zero bias.

\subsection{Training}

\label{sec:training}

To train a path scorer with stable supervision, each query must contribute equal gradient
signal regardless of its number of negative samples.
We construct a paired sampling
dataset that guarantees both balance and fairness.
\subsubsection{Paired Sampling}
\label{sec:paired-sampling}

For each query $q$, the query generation process produces one
\textit{positive} path $p^+_q$, which connects $s$ to the
true object $o$, and $N^-$ \textit{negative} paths
$\mathcal{S}^-(q) = \{p^-_{q,1}, \ldots, p^-_{q,N^-}\}$, which connect
$s$ to sampled false objects. The total number of pairs equals
$\sum_q |\mathcal{S}^-(q)|$, ensuring that every negative is visited exactly
once per epoch. Collation interleaves positive and negative pairs
into flat sequences
$[p^+_0, p^-_0, \ldots, p^+_{B-1}, p^-_{B-1}]$.
Here $B$ is the query batch size. Each training batch contains
exactly $B$ positive and $B$ negative samples, producing a guaranteed
balanced composition regardless of the underlying positive to negative
ratio. Variable length paths are zero-padded to the batch maximum
$N_{\max}(\mathcal{B}) = \max_{p \in \mathcal{B}} (\ell_p + 1)$
entity tokens, with padding masks $\mathbf{M}$ as defined in
Section~\ref{sec:token-interleaving}.

\subsubsection{Loss Function}
\label{sec:loss-functions}

PAMR is trained with a margin-based ranking loss that groups samples
by $(s, r, ?)$ and enforces a margin between the positive score
and all negative scores for that query. Let $\mathcal{B} = \{(\hat{y}_1, y_1), \ldots, (\hat{y}_{2B}, y_{2B})\}$
denote a training batch with $y_i \in \{0, 1\}$. The loss is defined as:
\begin{equation}
  \mathcal{L}_{\text{rank}} = \frac{1}{|\mathcal{Q}(\mathcal{B})|} \sum_{q \in \mathcal{Q}(\mathcal{B})}\; \frac{1}{|S^+(q)||S^-(q)|} \sum_{\hat{y}^+ \in S^+(q)} \sum_{\hat{y}^- \in S^-(q)}
  \max\bigl(0,\, m - \hat{y}^+ + \hat{y}^-\bigr)
  \label{eq:loss-rank}
\end{equation}
Here $\mathcal{Q}(\mathcal{B})$ denotes the set of unique queries in $\mathcal{B}$, and
$m = 1.0$ is the margin. This loss directly optimises the ranking objective and prevents
score collapse.

\section{Experiment}
\label{sec:experiment}
\subsection{Implementation Details}
\label{sec:implementation}

All experiments run on a Vast.ai instance equipped with an NVIDIA
GeForce RTX 5060 Ti and an AMD EPYC 7763 processor allocated with 8 threads
and 16\,GB virtual RAM. The software stack comprises Python~3.11,
PyTorch~2.7.1 with CUDA~12.8 and cuDNN~9.

\paragraph{Path extraction.}
The bidirectional BFS is implemented in C++ (GNU GCC 13.2, C++23) and called
from Python. The adjacency list is loaded once, queries are distributed over 8
threads, and per-thread outputs merge into one path table per split, cached on
disk as pickle files for reuse.

\paragraph{Training data and GPU memory.}
The paired dataset is rebuilt once per epoch from the precomputed path tables,
with per-query pairing across CPU cores; paths are zero-padded to the batch
maximum with the padding masks of Section~\ref{sec:paired-sampling}. The frozen
embedding table dominates GPU memory, which is managed through expandable
memory segments and periodic cache defragmentation after initialization.

\paragraph{Training configuration.}
PAMR is trained with AdamW~\cite{loshchilov2019decoupled} at learning
rate $\eta = 0.01$, weight decay $\lambda = 0.01$, $\beta_1 = 0.9$,
$\beta_2 = 0.999$. The ranking loss (Eq.~\ref{eq:loss-rank}) is
used with margin $m = 1.0$. Batch size is 8 queries (16 samples per
batch with the guaranteed balanced positive/negative composition).
Training proceeds for 50 epochs. The Transformer encoder uses
$N = 8$ layers and $h = 8$ attention heads with
$d' = 2048$ and dropout $p = 0.1$. The hidden dimension
$d = 768$ is auto-detected from the frozen VL-KGE
model (TransE, $D = 768$). The
VL-KGE base model remains frozen during the training of PAMR.

\subsection{Datasets}
\label{sec:datasets}

Experiments use WikiArt-MKG~v2~\cite{efthymiou2026vlkge}, a multimodal
knowledge graph of fine art with 224,166 entities (216,564 artworks and
7,602 attributes) and 22 relation types, split chronologically into
7,877,220 training, 208,513 validation, and 208,368 test triples.
Artworks carry image embeddings; styles, genres, years, and locations
carry textual ones, so entity types differ in what modalities they
expose. Training queries get one negative by rejection sampling;
validation and test queries get ten. Three relation types introduced as
preprocessing inverses
(\texttt{isRelatedToArtwork}, \texttt{isRelatedToArtist},
\texttt{hasCreatedArtwork}) are excluded from evaluation, leaving the
four artist-to-artist relations (\texttt{isInfluencedBy},
\texttt{isInfluencedOn}, \texttt{isPupilOf}, \texttt{isTeacherOf}).
These four connect artist entities, enabling evaluation under both
multimodal and non-multimodal reasoning, and they are the only
cross-modality relations in the dataset: query and answer differ in
modality, which is exactly when modality-as-entity path structure
carries multimodal information. To our knowledge WikiArt-MKG~v2,
published with VL-KGE, is the only public benchmark built this way. In
attachment-formulated graphs such as WN9-IMG or DB15K, each multimodal
feature binds to a single entity, so lifting it into a node would
produce a degree-1 leaf that mediates no path, and modality-as-entity
reasoning would be vacuous there. Porting PAMR to such graphs is
therefore a reformulation rather than a port, and we leave it to
future work.

Path extraction runs bidirectional BFS with at most $H = 5$ hops, base
fanout $F = 100$, and decay factor $\lambda = 0.5$
\cite{icaart26}. Both the full and the artist-only configurations are
evaluated, the concept-only configuration serving as the topology-only
counterpart described in Section~\ref{sec:path-extraction}.

\subsection{Evaluation Metrics}
\label{sec:evaluation-metrics}

For each query, the model assigns a scalar score to every candidate
object in the evaluated candidate set (the gold answer together with
its 10 sampled negative candidates, as described in
Section~\ref{sec:datasets}) and ranks them in descending order. We
report Mean Reciprocal Rank (MRR) and Hit-rate at cutoff $K$
(Hits@$K$), overall and per relation type.

\subsection{Results}
\label{sec:results}

Baselines come from VL-KGE~\cite{efthymiou2026vlkge} with four embedding
backbones: TransE, RotatE, DistMult, and ComplEx. LLM-based methods are
deliberately excluded since their pretraining over external corpora makes
any comparison unfair. Adapting message-passing path reasoners such as
NBFNet or RED-GNN to the modality-as-entity graph would require
redesigning their message functions around typed modality nodes; we leave
that adaptation to future work and use PAR, the concept-only path scorer,
as the structural baseline.

\subsubsection{Comparative Study}

Table~\ref{tab:results} reports macro-averaged MRR and Hits@$K$ across
the four evaluation relations between artists. PAMR and PAR operate over
the frozen TransE embeddings provided by VL-KGE ($D = 768$), the
translation-based, real-valued backbone that also produces the strongest
baseline; building on it asks whether path reasoning extracts more from
the same frozen representations than direct triple scoring does. The four
VL-KGE baseline columns come from
our own independent runs of the published VL-KGE scorers, one per
backbone, executed on the same query sets and sampling protocol as PAMR;
they are not the values printed in the VL-KGE paper, which aggregate over
the full relation set. The VL-KGE baselines score triples directly
without path reasoning, while PAR uses only structural paths and PAMR
additionally exploits multimodal token representations.

\begin{table*}[t]
  \centering
  \setlength{\tabcolsep}{2pt}
  \renewcommand{\arraystretch}{1.15}
  {\footnotesize
    \begin{tabular}{l|cccc|cc}
      \toprule
      \multirow{2}{*}{Metric} &
      \multicolumn{4}{|c|}{VL-KGE Embedding Baselines} &
      \multicolumn{2}{|c|}{Ours} \\
      \cmidrule(lr){2-5} \cmidrule(lr){6-7}
      & TransE & RotatE & DistMult & ComplEx & PAR & PAMR \\
      \midrule
      MRR  & \underline{0.363$\pm$0.030} & 0.327$\pm$0.029 & 0.311$\pm$0.020 & 0.320$\pm$0.048 & 0.221$\pm$0.027 & \textbf{0.379$\pm$0.042} \\
      Hits@1 & \textbf{0.387$\pm$0.044} & 0.182$\pm$0.069 & 0.318$\pm$0.076 & 0.283$\pm$0.061 & 0.058$\pm$0.027 & \underline{0.345$\pm$0.027} \\
      Hits@3 & \textbf{0.625$\pm$0.044} & 0.419$\pm$0.098 & 0.528$\pm$0.049 & 0.494$\pm$0.029 & 0.199$\pm$0.041 & \underline{0.576$\pm$0.060} \\
      Hits@5 & \underline{0.772$\pm$0.065} & 0.597$\pm$0.047 & 0.705$\pm$0.037 & 0.703$\pm$0.062 & 0.331$\pm$0.036 & \textbf{0.782$\pm$0.075} \\
      \bottomrule
  \end{tabular}}
  \caption{Macro-averaged ranking metrics across four evaluation
    relations between artists. \textbf{Bold} indicates the best per metric;
    \underline{underlines} indicate the second best. Values are reported
    as mean $\pm$ standard deviation. Baseline values are macro-averaged
    over the same four cross-modality relations from our independent runs
    and therefore do not correspond to the full relation set reported in
    the VL-KGE publication.}
  \label{tab:results}
\end{table*}

The ordering is consistent across method families: the four VL-KGE
embedding baselines occupy a moderate performance band, PAR falls
below it, and PAMR sits above. TransE leads the baselines with
0.363~MRR, while RotatE, DistMult, and ComplEx stay within
0.02~MRR of one another; among the three scoring functions beyond
TransE, the choice has limited impact on this graph.

Adding multimodal tokens carries PAMR past the strongest baseline:
0.379 MRR (a 4.4\% improvement over TransE) and the best Hits@5 at
0.782 (a 1.3\% improvement). At stricter cutoffs PAMR stays
competitive, 0.345 versus 0.387 at Hits@1 and 0.576 versus 0.625 at
Hits@3. The crossover, PAMR leading MRR and Hits@5 while TransE
leads Hits@1 and Hits@3, traces back to how the models spread
probability: TransE concentrates it on a few candidates, PAMR
distributes it, trading some top positions for depth of recall.
Computationally the trade is affordable; PAMR trains in about 543
seconds per epoch against TransE's 577 (training-step time,
excluding one-time path extraction and preprocessing,
Section~\ref{sec:runtime}). Per relation, PAMR leads most metrics
on \texttt{isPupilOf} and \texttt{isInfluencedOn}, TransE on
\texttt{isTeacherOf} and \texttt{isInfluencedBy}, and PAR trails on
all four.

\subsection{Runtime and Scalability Analysis}
\label{sec:runtime}

The pipeline splits into one-time offline preprocessing (graph ingestion,
query generation, bidirectional BFS extraction, and path-table
materialisation) and the recurring training loop.
Table~\ref{tab:runtime} reports the measured wall-clock time and peak
memory of each stage on the instance of
Section~\ref{sec:implementation}, single run, warm cache.

\begin{table}[t]
  \centering
  {\footnotesize
  \begin{tabular}{lrr}
    \toprule
    Stage & Wall-clock & Peak memory \\
    \midrule
    Graph ingestion (224{,}166 nodes) & \textbf{[TBD]}\,s & \textbf{[TBD]}\,MB \\
    Query generation & \textbf{[TBD]}\,s & \textbf{[TBD]}\,MB \\
    BFS, training split (7.87M queries) & \textbf{[TBD]}\,s & \textbf{[TBD]}\,MB \\
    BFS, validation split (208k queries) & \textbf{[TBD]}\,s & \textbf{[TBD]}\,MB \\
    BFS, test split (208k queries) & \textbf{[TBD]}\,s & \textbf{[TBD]}\,MB \\
    Path-table write & \textbf{[TBD]}\,s & \textbf{[TBD]}\,MB \\
    Paired-dataset build (per epoch) & \textbf{[TBD]}\,s & \textbf{[TBD]}\,MB \\
    Training epoch (PAMR) & 543\,s & \textbf{[TBD]}\,MB \\
    Training epoch (TransE baseline) & 577\,s & \textbf{[TBD]}\,MB \\
    Test inference (208k queries) & \textbf{[TBD]}\,s & \textbf{[TBD]}\,MB \\
    \bottomrule
  \end{tabular}}
  \caption{Measured per-stage wall-clock time and peak memory on the
    instance described in Section~\ref{sec:implementation}.}
  \label{tab:runtime}
\end{table}

Preprocessing dominates the offline budget: with 7.87M training queries,
bidirectional BFS over the full graph takes \textbf{[TBD]}\,s, while the
much smaller validation and test splits finish in \textbf{[TBD]}\,s and
\textbf{[TBD]}\,s. The offline cost is paid once per experiment; the
extracted paths are cached as path tables and reused across training
runs, so the 543\,s/epoch of PAMR against TransE's 577\,s/epoch reflects
the recurring training cost alone.

The sensitivity of the extraction stage to its two hyperparameters, the
hop bound $H$ and the fanout $F$, is summarised in
Table~\ref{tab:bfs-sweep}. At the operating point $H = 5$, $F = 100$
used throughout the paper, BFS completes in \textbf{[TBD]}\,s with
\textbf{[TBD]}\,MB of peak memory. Reducing the hop bound to
$H = \textbf{[TBD]}$ lowers the cost to \textbf{[TBD]}\,s but caps the
reachable path length at $L = 2H + 3$ tokens, and raising the fanout to
$F = \textbf{[TBD]}$ increases it to \textbf{[TBD]}\,s; the frontier
decay keeps the growth well below the naive $O(F^{2H})$ bound in
practice. The chosen operating point therefore caps token sequences at
$L = 13$ while keeping extraction within \textbf{[TBD]}\,s of the
cheapest configuration measured.

\begin{table}[t]
  \centering
  {\footnotesize
  \begin{tabular}{lrr}
    \toprule
    Configuration & Wall-clock & Peak memory \\
    \midrule
    $H = 2$, $F = 100$ & \textbf{[TBD]}\,s & \textbf{[TBD]}\,MB \\
    $H = 3$, $F = 100$ & \textbf{[TBD]}\,s & \textbf{[TBD]}\,MB \\
    $H = 4$, $F = 100$ & \textbf{[TBD]}\,s & \textbf{[TBD]}\,MB \\
    $H = 5$, $F = 100$ (default) & \textbf{[TBD]}\,s & \textbf{[TBD]}\,MB \\
    $H = 5$, $F = 50$ & \textbf{[TBD]}\,s & \textbf{[TBD]}\,MB \\
    $H = 5$, $F = 200$ & \textbf{[TBD]}\,s & \textbf{[TBD]}\,MB \\
    \bottomrule
  \end{tabular}}
  \caption{BFS extraction wall-clock time and peak memory as the hop
    bound $H$ and fanout $F$ vary (training split, 8 threads).}
  \label{tab:bfs-sweep}
\end{table}

On the training side, scoring a candidate costs $O(L^{2}h)$ attention
over its token sequence, and per-query inference latency grows with the
candidate set size. GPU memory is dominated by the frozen embedding
table; the trainable Transformer is small. These figures locate the
practical limits of the design: preprocessing scales with query volume
and graph density, scoring scales with the candidate pool, and both
bottlenecks are revisited in Section~\ref{sec:conclusion}.

\subsection{Case Study}
\label{sec:case-study}

Figure~\ref{fig:case-study} shows where multimodal entities matter
most. The query (Pablo Picasso, \texttt{isInfluencedOn}, ?) is taken
verbatim from the test set, with two candidates to compare: Salvador
Dali and Lee Krasner.

\begin{figure}[t]
  \centering
  \includesvg[width=\linewidth]{figures/case.svg}
  \caption{Case study on the test query (Pablo Picasso,
    \texttt{isInfluencedOn}, ?): candidates Salvador Dali and
    Lee Krasner, a subgraph of intermediate artist relations, and
    the artworks \textit{Man with a Porron} (by Dali),
    \textit{Child with a Dove} (by Picasso), and \textit{Self-Portrait}
    (by Matisse) as multimodal nodes.}
  \label{fig:case-study}
\end{figure}

Restricted to artists, PAR's bidirectional BFS can only walk through colleague relations: Picasso $\xrightarrow{\text{isRelatedToArtist}}$ Juan Gris $\xrightarrow{\text{isRelatedToArtist}}$ Dali, and Picasso $\xrightarrow{\text{isRelatedToArtist}}$ Georges Braque $\xrightarrow{\text{isRelatedToArtist}}$ Matisse $\xleftarrow{\text{isInfluencedBy}}$ Krasner. Both are plausible social routes, but neither carries anything artistic, so the two candidates are hard to tell apart.

With artworks as nodes, PAMR's BFS routes through the creative works. The Dali path re-enters via Picasso's \textit{Child with a Dove} $\xrightarrow{\text{isRelatedToArtwork}}$ \textit{Man with a Porron} $\xrightarrow{\text{hasArtist}}$ Dali; the Krasner path runs through Picasso $\xleftarrow{\text{hasArtist}}$ \textit{Child with a Dove} $\xrightarrow{\text{isRelatedToArtwork}}$ \textit{Self-Portrait} $\xrightarrow{\text{hasArtist}}$ Matisse $\xleftarrow{\text{isInfluencedBy}}$ Krasner.

Given these token sequences, PAMR ranks Dali above Krasner. When the frozen artwork embeddings pass through the encoder, they carry the visual and semantic kinship of the pieces along with them. The Dali route keeps to Spanish avant-garde territory, with shared structural experimentation and foundational Cubist motifs linking Picasso's \textit{Child with a Dove} and Dali's \textit{Man with a Porron}. Reaching Krasner through Matisse's \textit{Self-Portrait}, by contrast, demands a leap from European Cubism through Fauvism to American Abstract Expressionism, movements separated in style, geography, and time, and the path is scored accordingly.

Although the Transformer scorer itself is a learned and non-transparent component, its judgement here rests on paths a reader can audit: which artworks the route visits, and what each visit contributes. Compared with the purely conceptual graph, the reasoning trajectory is grounded in art-historical facts.

\section{Conclusion}
\label{sec:conclusion}

This paper introduced PAMR, a path-aware reasoner for multimodal knowledge graphs that operates under a modality-as-entity setting, multimodal values become graph nodes, and structural paths are scored by a Transformer over frozen KGE tokens. The framework attained the best MRR on WikiArt-MKG v2's four cross-modality artist relations and, more importantly for multimodal practice, attaches to each prediction a path through shared artworks and attributes that a reader can follow end to end. The learned scorer itself remains non-transparent, but what it scores is auditable.

The study has limits. The scoring spreads probability widely, which serves MRR and Hits@5 but loses Hits@1 and Hits@3 to TransE, possibly because eight encoder layers are not enough to sharpen the head of the ranking; inference over long interleaved sequences also stays costly when graphs grow. Future work will target both: KV caching and hardware-aware attention for long sequences, deeper encoders for the top ranks, and contrastive objectives that separate top candidates from hard negatives. Extending the modality-as-entity formulation to dynamic or temporal modalities is the natural next benchmark, and, as traceable evidence becomes a central requirement of trusted multimodal systems, we consider path-grounded judgement a strength worth keeping in future designs.

\begin{credits}
  \subsubsection{\ackname} This research is funded by the University of Science, VNU-HCM, Vietnam under grant number CNTT 2025-16. The authors hereby disclose that an AI tool, namely Google Gemini, has been used to perform literature reviews as well as language and text improvements, and the authors declare that all of the texts generated by the AI have been carefully reviewed manually. No illustration was made with an involvement of any AI tool.

  \subsubsection{\discintname}
  The authors have no competing interests to declare that are
  relevant to the content of this article.
\end{credits}

\bibliographystyle{splncs04}
\bibliography{refs}

\end{document}
